% LaTeX rendering of paper5-authority-derivation-draft-v0.6-populated.md
% (the consolidated manuscript). Hand-typeset, modeled on sarc-suite-one-pass's
% own paper-tex/main.tex conventions (commit 782261e, ADR-001-foundation.md) --
% same document class, packages, and theorem-environment setup, content
% written fresh for this paper. This repo ports only G1 (build, byte-
% identical double build under a content-stable epoch), G7 (arXiv sidecar
% sync), G8 (citation completeness), and G9 (bibliography quality) from the
% sibling's own G1-G9 gate suite -- not G2-G6 (the sibling's own tex-vs-
% markdown token/number/prose/structure parity gates), so this file is a
% faithful rendering of the same real, already-verified figures, not a
% mechanically parity-checked transliteration.
%
% Canonical build toolchain: pinned Tectonic 0.17.0 (bootstrap.sh, via
% sarc-suite-one-pass's own bootstrap.sh delegation, ADR-001-foundation.md).
% Build manually: `tectonic main.tex`; `make arxiv` (this repo's own
% Makefile) builds and runs gates/run_gates.py together.
\documentclass[11pt,a4paper]{article}

\usepackage[utf8]{inputenc}
\usepackage[T1]{fontenc}
\usepackage{lmodern}
\usepackage{amsmath}
\usepackage{amssymb}
\usepackage{amsthm}
\usepackage{booktabs}
\usepackage{longtable}
\usepackage{array}
\usepackage{microtype}
% Unlike the sibling's own main.tex, this file has no G4 word-for-word
% prose-parity gate to preserve (task brief: only G1/G7/G8/G9 are
% ported), so ordinary dictionary hyphenation is left enabled -- it is
% what actually resolves an overfull \hbox on a long unbreakable
% identifier or URL, rather than working around one with extra stretch.
\usepackage[margin=1in]{geometry}
\usepackage[colorlinks=true,linkcolor=blue,citecolor=blue,urlcolor=blue]{hyperref}
\setlength{\emergencystretch}{8em}

\theoremstyle{definition}
\newtheorem{definition}{Definition}

\newcommand{\proofstatus}[1]{[\textsc{PROOF-STATUS}: \textsc{#1}]}
\newcommand{\pstar}{P\textsuperscript{*}}

\title{Minimal Sufficient Governance Context: Compiling Loss Models and Reachable Execution Semantics into Cost-Aware Runtime Authority Contracts for Autonomous Agents}
\author{Gaston Besanson\thanks{Universidad Torcuato Di Tella}}
\date{}

\begin{document}
\maketitle

\noindent Companion artifact: \texttt{sarc-authority-derivation}. Paper 5 of the SARC series, built on the pinned \texttt{sarc-suite-one-pass} artifact (arXiv 2608.18360, commit \texttt{782261e}) as a read-only imported baseline.

\begin{abstract}
Given a declared loss model and the states an autonomous system can reach, we derive and certify the minimal, cost-aware sets of runtime observations sufficient to make every modeled safe-versus-unsafe authority distinction. We do not declare an observation set and hope it is complete: we \emph{synthesize} every minimal sufficient contract (a reduct, encoding this artifact's own discernibility sets -- Skowron and Rauszer, 1992 -- as CNF for SAT/MaxSAT solvers), certify each one with a machine-checkable partition certificate, and select among them by declared observation cost, not cardinality alone. On a real code/cloud authority domain (not planted for this purpose), the individually-indispensable core is \textbf{not sufficient} as a runtime contract (six of ten candidates) and two distinct seven-attribute reducts exist; a preregistered cost model separates them exactly (CH-B2: \textbf{supported}, the cheaper contract costs 6.124 against 7.354 for the alternative, both independently certified safety-equivalent). On a second, larger, non-planted domain built independently of this question (thirty-five candidate properties), the same core-is-not-a-reduct distinction recurs (CH-C1: \textbf{supported}, a four-property core, minimum reduct cardinality nine, at least five distinct minimum-cardinality contracts) -- but the same domain's own cost model does \emph{not} separate them: CH-C2 is \textbf{not supported}, a genuine, fully explained cost tie across every known minimum-cardinality contract at 9.153, reported as found, not reframed. We register and measure discernibility-family scaling on candidate-attribute universes up to one hundred properties, where exhaustive reduct enumeration is expected and confirmed infeasible within a registered three-hundred-second budget, while SAT/MaxSAT synthesis solves in well under a second; MaxSAT showed no measured cardinality advantage over plain SAT on these specific registered families, a two-sided registered question, not a claim tuned toward the outcome we expected. AuthorityBench compares four baselines -- a declared-only manual policy, ABAC-mining reimplemented from Xu and Stoller (2015), essential-variable analysis, and exhaustive reduct enumeration -- across three domains; the declared-only baseline is not exactly sufficient on any of them. Every synthesized contract is compiled with a sufficiency certificate and, where the reachable set is enumerable, a counterexample when it is not -- a runtime gate is meant to consume the certificate, not a claim.
\end{abstract}

\section{Problem}

A pre-action authority gate is exactly as complete as whoever declared its observation set thought to be. \texttt{sarc-suite-one-pass}'s own imported baseline (\texttt{specs/authority.yaml}) observes two things: whether the acting role is on an allow-list, and whether an order value exceeds a single scenario-wide cap -- and, on the real code/cloud domain below, a \textbf{declared-only policy of this same shape is not exactly sufficient} (a stronger, machine-checked version of the same finding this artifact's own v2 baseline comparison already made, and a finding that also holds on AuthorityBench's other two registered domains, Section~\ref{sec:authoritybench}). This paper's claim: given one declared loss model and one executable-reachable action set, \emph{compile} -- do not hand-author -- the exact, minimal, cost-aware runtime context a gate must observe, with alternative contracts, certificates, and counterexamples a downstream system can consume directly, not merely a number a human reads once.

\section{Formal Objects}
\label{sec:formal}

Let a \textbf{reachable state} be one executable-reachable tuple over a declared set of candidate properties. Let a \textbf{loss model} $M$ be a finite set of boolean predicates over a reachable tuple, each declaring one way an action can be unsafe. Two reachable tuples are a \textbf{cross-verdict pair} if $M$'s verdict differs between them.

\begin{definition}[Participation, the one-flip criterion]
Property $p$ is authority-bearing for $M$ iff there exist executable-reachable tuples $a, a'$ differing only in $p$ -- a single flip, holding every other candidate property fixed -- whose $M$-verdicts differ.
\end{definition}

\begin{definition}[The core \pstar]
\pstar\ is the set of every property participating by Definition~1.
\end{definition}

\begin{definition}[The derived coverage list]
The candidate properties not in \pstar.
\end{definition}

\begin{definition}[Sufficiency]
A property set $S$ is sufficient for $M$ iff any two reachable tuples agreeing on every property in $S$ also agree on their $M$-verdict. Equivalently (Skowron and Rauszer's discernibility-matrix device, 1992): for every cross-verdict pair, its \emph{difference set} is the candidate properties on which the pair disagrees, and $S$ is sufficient iff $S$ hits every difference set.
\end{definition}

\begin{definition}[A reduct]
A minimal sufficient set: sufficient, and no proper subset is.
\end{definition}

\begin{definition}[The core, restated as an intersection]
The intersection of every reduct.
\end{definition}

This discernibility restatement is what makes reduct synthesis tractable past exhaustive enumeration (Section~\ref{sec:scale}): a difference set that is a proper superset of another already in the family adds no separate hitting-set constraint, and the pruned family becomes one CNF clause per surviving difference set and one boolean variable per candidate property -- a SAT or MaxSAT solver's own input, not the $2^n$ subset enumeration every backend is still cross-checked against everywhere exhaustion stays affordable.

\section{Core versus Reduct}
\label{sec:core-vs-reduct}

Definition~2's \pstar\ (individual indispensability, a singleton-perturbation witness) and Definition~5's reduct (joint sufficiency, Definition~4) are different properties.

\noindent\textbf{Proposition 1$'$ (replacing the original, corrected Proposition~1).} Claim one, Definition~1's witness search computes exactly the core, Definition~6; claim two, the core is a subset of every reduct; claim three, \emph{if} the core is sufficient, \emph{then} the core is the unique reduct -- claim three's antecedent never assumed, always decided per instance by a sufficiency certificate or a concrete counterexample.

\noindent\textbf{Negative Proposition N (the core need not be sufficient, in general).} The registered two-state counterexample (candidate properties $x, y \in \{0,1\}$, reachable set $\{(x{=}0,y{=}0), (x{=}1,y{=}1)\}$, verdict $M(x,y) = x$) has an \emph{empty} core -- no reachable pair differs in exactly $x$ or exactly $y$, since the set's only two members differ in both at once -- that is trivially \emph{not sufficient} ($M(0,0) \neq M(1,1)$, yet the empty set determines nothing), with $\{x\}$ and $\{y\}$ each an independent singleton reduct, neither equal to the core. This is the sharpest possible instance: a core that is not merely non-maximal but \emph{empty}, directly falsifying the general claim the original Proposition~1 made (Section~\ref{sec:corrections} records the correction itself as a finding, not smoothed over).

\noindent\textbf{The worked introduction: this artifact's own v2 procurement instance, where claim three's antecedent holds.} This artifact's registered retail-procurement domain -- nine original candidate properties plus \texttt{min\_order\_quantity}, a declared cross-field reachability filter, and two characterized remediation mechanisms -- is not a counterexample to Proposition~1$'$: it is the case where the antecedent fires. CH-A8 certifies the v2 core sufficient over the full reachable set; CH-A10 (exhaustive enumeration of every candidate subset) confirms the core -- nine properties -- is the \textbf{unique reduct}, of minimum cardinality nine. This is this paper's simplest instance, worked first because it is the one case in this whole artifact where core and reduct coincide -- every domain in Section~\ref{sec:evidence} is chosen to test whether that coincidence holds again, and mostly does not.

\section{Realistic Evidence: The Empirical Centre}
\label{sec:evidence}

Section~\ref{sec:core-vs-reduct}'s abstract distinction and its one worked coincidence (v2) leave open whether either shape is typical on a domain not designed around the question. Four hypotheses, on two independently-built domains neither planted to exhibit core-versus-reduct or cost separation, are this paper's own empirical centre.

\noindent\textbf{CH-B1 (a software-and-cloud execution-agent domain).} \proofstatus{machine-checked} (\texttt{checkers/ch\_b1\_check.py}, all candidate subsets of ten candidates considered, subset-pruned). The core (six properties: \texttt{approval\_token, data\_classification, delegated\_role, operation, repository, resource\_owner}) is \textbf{not sufficient} as a contract -- individually indispensable, not jointly enough -- and two distinct minimum reducts exist, both of cardinality seven, each the core plus exactly one more property: the core plus \texttt{branch}, or the core plus \texttt{environment} -- two operationally different, equally sufficient ways to close the same coverage gap.

\noindent\textbf{CH-B2 (does a preregistered observation-cost model separate CH-B1's own two reducts?)} \proofstatus{machine-checked} (\texttt{checkers/ch\_b2\_check.py}). \textbf{Supported}, positive strict cost separation: the minimum-cost contract (core plus \texttt{branch}, cost 6.124) is strictly cheaper than the other minimum-cardinality reduct (cost 7.354, delta 1.230) -- both independently certified safety-equivalent before cost is consulted at all, so cost breaks a tie between two contracts already known interchangeable for safety, not a shortcut around safety.

\noindent\textbf{CH-C1 (does core-versus-reduct recur on a second, independently-built, much larger domain?)} A thirty-five-property enterprise access-governance domain -- seven properties from each of five families, six reachability drivers, twenty-nine derivation rules -- whose reachability and derivation structure were fixed for organisational plausibility \emph{before} any core, reduct, or cost outcome was computed. \proofstatus{machine-checked} (\texttt{checkers/ch\_c1\_check.py}, 27{,}000 reachable tuples swept). \textbf{Supported}: the core (four properties: \texttt{approval\_token, delegated\_role, resource\_environment, workflow\_stage}) is \textbf{not sufficient}, minimum contract cardinality is nine, and at least five distinct minimum-cardinality contracts exist -- a genuine lower bound, via blocking-clause enumeration capped at five, never claimed exhaustive. Exhaustive enumeration was separately attempted, registered budget three hundred seconds, and confirmed infeasible (wall time 300.08s, against 6{,}724{,}520 size-seven subsets alone).

\noindent\textbf{CH-C2 (does a preregistered cost model separate CH-C1's own minimum-cardinality contracts?)} \proofstatus{machine-checked} (\texttt{checkers/ch\_c2\_check.py}, same domain as CH-C1). \textbf{Not supported}, a negative tie: every one of the five known minimum-cardinality contracts costs exactly 9.153 under this domain's own seven-tier cost model -- a genuine, fully explained tie, retained and reported, not reframed. The tied contracts substitute properties from the \emph{same} declared cost tier for one another, so a cost model built at this granularity cannot discriminate between them -- weighted optimisation over declared cost is \textbf{conditional} on the cost model actually separating the substitute properties in play; CH-B2 shows a domain where it does, CH-C2 shows one where, at the granularity actually declared, it does not.

\section{Synthesis: SAT, MaxSAT, and Weighted Contracts}
\label{sec:synthesis}

\texttt{synthesis.py} encodes Section~\ref{sec:formal}'s discernibility family as CNF and exposes three backends: plain SAT (any sufficient contract, not necessarily minimal), cardinality MaxSAT (a soft unit clause per property, weight one), and weighted MaxSAT (the same construction with each property's declared observation cost as its own soft-clause weight); a blocking-clause multiplicity method (added for CH-C1) repeatedly re-solves the same objective, excluding each prior model with one added hard clause, stopping as soon as an iteration's optimum exceeds the first -- a genuine lower bound on multiplicity, never $2^n$ enumeration. Every backend is held to exhaustive reduct enumeration's own from-first-principles answer everywhere exhaustion is affordable -- a real regression would be caught by comparison, not assumed away by trusting a solver.

Cost-aware synthesis' mechanism was proved first on a small, hand-derived domain before it was asked of a realistic one: a boolean \texttt{a} and a bijected pair \texttt{b1, b2} where minimum cardinality (\texttt{\{a\}}, cost 4.00) and minimum declared cost (\texttt{\{b1, b2\}}, cost 1.00, delta 3.00) disagree by construction -- proof the mechanism \emph{can} select a strictly cheaper sufficient contract over a smaller one, not yet evidence a realistic domain does the same. This artifact's own v1/v2/v4 domains showed no divergence at the time (uniform declared cost floor, every one of their own cost deltas exactly zero) -- the operational-value question this left open is exactly what CH-B2 and CH-C2 above now answer, on two realistic domains, in both directions.

\section{Scale: Symbolic Synthesis Beyond Exhaustive Search}
\label{sec:scale}

Three registered scaling results, kept as three separate, still-valid questions, none superseding another.

\noindent\textbf{Exploratory (v5.0, kept, degeneracy disclosed).} The original planted-reduct scaling experiment grew the candidate-property count alone ($n = 10, 30, 50, 100$) while deliberately holding the reachable set fixed at six tuples. All exact at every size, wall time between 0.0004s and 0.0005s with no distinguishable growth trend -- a real, valid answer to the question it was built for, but a \textbf{degenerate} one for scale in general: the noise properties never appear in any discernibility-family clause, costing the solver nothing beyond one extra boolean variable each. Relabelled \texttt{exploratory\_v5\_0} and superseded as the primary scaling claim by the two results below, not deleted.

\noindent\textbf{The tuple-scaling result (v5.1, kept).} Holding the candidate-property count small enough to stay exhaustible, a planted two-reduct family grows the reachable set itself to $100{,}001$ tuples: the discernibility family still reduces, after superset removal, to exactly one set, and both planted reducts are recovered exactly by every backend.

\noindent\textbf{The primary scaling table (v5.3): combinatorial hardness.} Three families pair a planted minimum-reduct size with a large candidate-attribute universe, far past exhaustive-enumeration tractability, and a diverse discernibility structure:

\begin{center}
\begin{tabular}{lcccc}
\toprule
Family & $n$ & $k$ & Discern.\ family (pruned) & Exhaustive enumeration \\
\midrule
A & 30 & 5 & 125 sets & infeasible, 300s budget \\
B & 60 & 10 & 250 sets & infeasible, 300s budget \\
C & 100 & 15 & 375 sets & infeasible, 300s budget \\
\bottomrule
\end{tabular}
\end{center}

Cardinality- and cost-MaxSAT both found the exact, hand-proved minimum reduct at every family regardless, in well under a second per family -- families that are hard for exhaustive enumeration are not automatically hard for the solvers synthesizing over their own discernibility family (Section~\ref{sec:limitations}). \textbf{Two-sided, reported as measured}: MaxSAT showed no meaningful cardinality advantage over plain SAT on any of these three families -- this registration pre-committed to accept and report a uniformly-absent advantage as a legitimate outcome, not a benchmark failure.

\section{AuthorityBench}
\label{sec:authoritybench}

AuthorityBench runs four baselines -- manual least-privilege (the declared-only shape), ABAC policy mining (reimplemented from Xu and Stoller, TDSC 2015, independently validated against the published ``Health Care Sample Policy'' case study before being trusted here), essential-variable analysis, and exhaustive reduct enumeration -- across three domains this artifact independently registered (procurement, code/cloud, data-and-communications), six metrics each. The declared-only manual baseline is not exactly sufficient on any of the three -- Section~1's opening claim, measured across every registered domain this artifact has, not asserted from one. Mining agrees with the packaged \texttt{derive\_authority\_contract}'s own minimum-cardinality answer on the procurement and code/cloud-core instances and lands on the \emph{other} known minimum reduct (the \texttt{environment} variant of CH-B1's own two) on the code/cloud domain -- both genuine, independently-verified minimum reducts of the identical model, a tie-break divergence between two independently-computed global minima, not a disagreement about which properties matter. Decision: \textbf{inconclusive} by the registered rule applied honestly to what was measured -- a registered non-result for whether mining tends to be avoidably non-minimal, reported as such rather than reframed as a win.

Two further registered results share this domain portfolio. \textbf{The budget-binding scenario (v5.2)}: a procurement scenario constructed so the budget loss predicate actually fires -- it did, in every one of sixty of sixty cells (137{,}711 firings total) -- and CH-A1's own derived-zero-miss guarantee held on every one of them. \textbf{The isolated-arm re-measurement (v0.4)}: direct code inspection found the baseline, derived, and over-inclusive experiment arms sharing one budget and one grant ledger, a real internal-validity threat registered before any fix was written. The isolation hypothesis itself was \textbf{not supported}: all thirty registered seeds across both workflows, re-run under the corrected, isolated design, reproduce the prior shared-state result's own means, confidence intervals, and booleans byte-for-byte -- a direct diagnostic confirmed the one loss predicate the bug could have corrupted fired zero times in any cell. The corrected design is kept regardless of the null result.

\section{Runtime Implication}

\texttt{src/authority\_compiler} packages this paper's own machinery behind one call, \texttt{derive\_authority\_contract}, returning a compiled \texttt{AuthorityContract}: the core, every reduct where exhaustion is affordable, the minimum-cardinality and minimum-cost reducts, the redundant-attribute list, a sufficiency certificate (the partition itself, not a bare boolean), counterexamples when the core is not sufficient, and a bound \texttt{contract\_change\_delta} callable reporting whether a newly reachable set of tuples still leaves a prior contract sufficient before resynthesizing anything. A runtime pre-action gate is meant to consume this compiled contract and its certificate directly -- gate schema and enforcement code are downstream engineering this paper does not implement.

\section{Related Work and the Novelty Fence}

Nine neighbouring literature clusters, each fetch-verified: what each established, and what this paper adds.

\textbf{STPA and its descendants} \cite{stpa-handbook,stpa-sec,phase-rismani-2024,mylius-2025,deepstpa} establish a structured but human-driven analyst process: an engineer enumerates losses and hazards and writes safety constraints by hand, checked by expert review. None derive, in general, a machine-checked minimal property set for a runtime gate, or formalize soundness or minimality as machine-checkable claims over an enumerated model. This paper's derivation is inspired by STPA's loss-to-constraint direction of travel but replaces the analyst step with a deterministic, exhaustively-verified derivation.

\textbf{ABAC policy mining} \cite{xu-stoller-abac-mining,abac-mining-survey-2022} starts from an existing lower-level policy plus attribute data and reconstructs which attributes already participate. This paper runs the derivation in the opposite direction, from a declared loss model forward to the minimal property set a new gate must observe, with a machine-checked soundness and minimality guarantee that mining, in general, has no ground-truth policy to verify against and so cannot provide.

\textbf{Policy-language completeness} \cite{ptacl,sacmat16-completeness} establishes what a policy combination language can express once the participating attributes and base decisions are already fixed; neither asks which attributes must be observed in the first place. This paper contributes exactly the question both assume already answered.

\textbf{Non-interference} \cite{goguen-meseguer-1982} is a semantic criterion for whether one party's actions can affect what another observes -- a check of dependence between two fixed parties, in general neither a derivation procedure nor a minimality guarantee. Definition~1 borrows the underlying comparison but turns it into a derivation swept over every candidate property against every loss predicate, with soundness and minimality verified exhaustively.

\textbf{Counterfactual fairness} \cite{counterfactual-fairness} is structurally the same comparison Definition~1 uses, applied once, to a single pre-chosen sensitive attribute. This paper generalizes the comparison to every candidate property against a declared loss model and adds an exhaustive, machine-checked minimality result that, in general, the counterfactual-fairness literature neither claims nor checks.

\textbf{Shield synthesis} \cite{bloem-shield-synthesis-2015} synthesizes a runtime monitor correcting violating outputs over a system's already-declared interface signals. This paper solves the strictly prior problem, in general, with a machine-checked minimality guarantee, upstream of the corrective strategy shield synthesis solves for.

\textbf{Runtime enforcement and edit automata} \cite{schneider-enforceable-security-policies-2000,ligatti-bauer-walker-edit-automata-2005} characterize and extend what a monitor may do to an already-fixed target-action alphabet, never which properties must be part of that alphabet in the first place -- the antecedent question this paper answers in general, with a minimality guarantee neither work needs for its own narrower question.

\textbf{Controller synthesis} \cite{ramadge-wonham-supervisory-control-1987,pnueli-rosner-reactive-synthesis-1989} presupposes the plant's or module's own observable/controllable alphabet and synthesizes a control law within it. In general, the classical minimality result here is about the supervisor's behaviour given that alphabet, not the alphabet itself; this paper derives the alphabet a controller must observe, before any supervisor or reactive module is synthesized over it.

\textbf{Capability systems} \cite{dennis-van-horn-capabilities-1966} establish the capability as an unforgeable reference carrying authority -- \emph{how} authority, once granted, is represented and delegated, not which conditions must be checked before it is minted. This artifact's own execution-grant mechanism is engineering informed by this lineage's modern descendants \cite{macaroons-ndss-2014,ucan-specification,biscuit-specification}, not a novelty claim in its own right.

\noindent\textbf{The fence, stated once.} In general, none of the nine clusters above derive, from a declared loss model, a machine-checked minimal property set for a runtime authority gate, with soundness and minimality verified exhaustively over an enumerated model and certificates a downstream system consumes directly. This paper's contribution sits upstream of all nine, answering a question each of them assumes already settled.

\section{Limitations}
\label{sec:limitations}

\begin{itemize}
\item \textbf{Declared loss-model completeness} is assumed, not verified: a declared model missing a real hazard still produces a certified core or reduct \emph{for the declared model}, silently uninformative about the missing one.
\item \textbf{Reachability-model validity}: every result above is exact \emph{given} its own declared reachable set, transcribed from its own registration, not measured from a live deployment.
\item \textbf{Partial observability} is not modeled: every synthesis backend assumes the candidate properties it is given are the ones a gate \emph{could} observe.
\item \textbf{Temporal and probabilistic losses} are out of scope: every loss predicate here is a boolean function of one reachable tuple.
\item \textbf{External validation}: the reproduction packet is this artifact's own attempt to make independent reproduction possible; per its own eight-point standard, it is not itself independent evidence until someone who did not build this artifact runs it and reports back (no reproduction report filed as of this draft).
\item \textbf{Families hard for exhaustive enumeration are not automatically hard for the solvers} synthesizing over their own discernibility family: Section~\ref{sec:scale}'s own v5.3 families are registered-infeasible to enumerate and solve in well under a second regardless.
\item \textbf{Reduct multiplicity in the Section~\ref{sec:evidence} large domain (CH-C1) is a lower bound}, not an exhaustive count: blocking-clause enumeration capped at five found five distinct minimum-cardinality contracts; whether more exist above the cap is open.
\end{itemize}

\section{Corrections and Historical Foundation}
\label{sec:corrections}

Append-only: nothing below is retracted or reworded; each correction is recorded alongside its original claim.

\noindent\textbf{The v1 singleton-sufficiency overclaim (corrected v0.3).} The original Proposition~1 asserted that observing exactly \pstar\ lets a gate always compute $M$'s true verdict on any executable-reachable tuple -- conflating individual indispensability with joint sufficiency. A commissioned external, AI-assisted research review identified the mischaracterization; per this project's own discipline the finding was not applied directly -- it was adjudicated against this artifact's own definitions, and Negative Proposition N's counterexample was re-derived from first principles. Proposition~1 is superseded by Proposition~1$'$ (Section~\ref{sec:core-vs-reduct}), not deleted.

\noindent\textbf{The v5.0 degenerate scaling benchmark (disclosed v5.1).} Section~\ref{sec:scale}'s own exploratory benchmark scaled candidate-property count while holding the reachable set fixed at six tuples -- a valid answer to the question it was built for, but one where the added noise properties cost the solver nothing beyond one boolean variable each. Relabelled \texttt{exploratory\_v5\_0} from v5.1 forward, superseded as the primary scaling claim, kept and cited.

\noindent\textbf{The entangled v0.3 experiment arms and their null delta (found and corrected v0.4).} Section~\ref{sec:authoritybench}'s own budget/ledger-sharing defect, found by this paper's own code inspection, registered two-sided before any fix was written: isolating the three experiment arms' state produced a result byte-identical to the prior shared-state measurement -- \textbf{not supported}, with the mechanism explaining why rather than left unexplained.

\noindent\textbf{The two commissioned external review rounds.} Round one produced the core-versus-reduct correction above. Round two produced the broadened contribution this paper's own front matter names -- the code/cloud flagship promoted to the empirical centre, packaging and CI-smoke-test gates, symbolic scaling and AuthorityBench framing, and the large non-planted domain this draft adds. Neither round's finding was ever applied as a direct patch: every finding this paper's own results reflect was independently re-derived and machine-checked against this artifact's own definitions and data before being accepted.

\noindent\textbf{Where the full historical procurement paper lives.} Versions 0.1 through 0.5 are each frozen at their own commit, unedited since, and linked from the companion repository's own \texttt{README.md} History section -- the complete v1-through-v4 Introduction, novelty fence, replication, Definitions~1--6, derivation procedure, grant mechanism, pair testing, empirical section, and their own Threats to validity, Limitations, and Conclusion are preserved there in full, not reproduced a second time in this consolidated manuscript.

\section{Threats to Validity, Claims, and Proof Status}

\noindent\textbf{Threats to validity, written adversarially.} First, every ``realistic'' domain in Section~\ref{sec:evidence} is still a declared model, not a live deployment -- a domain built by the same author to be plausible is not independent evidence that real deployments share its structure. Second, CH-C1's own multiplicity count is a lower bound capped at five by construction -- a reader could reasonably suspect the true count is far larger, which would weaken, not strengthen, any claim resting on ``few'' alternatives. Third, CH-C2's cost tie could be read as this paper quietly avoiding a more impressive result -- the tie is reported at face value, with the mechanism that produces it stated explicitly. Fourth, AuthorityBench's mining comparison is inconclusive by its own registered rule; it is reported because the registration committed to reporting whatever was measured, not only a clean win.

\begin{center}
\begin{longtable}{p{0.47\textwidth}p{0.47\textwidth}}
\toprule
This paper claims & This paper does not claim \\
\midrule
\endhead
On two independently-built, non-planted domains, the core is not always a reduct (CH-B1, CH-C1). & That the core is never a reduct -- v2's own worked instance (Section~\ref{sec:core-vs-reduct}) is a real counterexample to that stronger claim. \\
On one of those two domains, declared cost strictly separates the minimum-cardinality alternatives (CH-B2). & That declared cost always separates them -- CH-C2, on the other domain, is a genuine tie. \\
Exhaustive reduct enumeration is registered-infeasible on the Section~\ref{sec:scale} hardness families and the Section~\ref{sec:evidence} large domain, and SAT/MaxSAT synthesis solves them regardless. & That every hard-to-enumerate instance is easy for these solvers, or that this artifact has found the general boundary between the two. \\
MaxSAT showed no measured cardinality advantage over plain SAT on the specific registered v5.3 families. & That MaxSAT never has a cardinality advantage over SAT in general. \\
The declared-only manual baseline is not exactly sufficient on any of AuthorityBench's three registered domains. & That no declared-only policy could ever be sufficient -- v2's own core happens to be sufficient (Section~\ref{sec:core-vs-reduct}). \\
The reproduction infrastructure is complete and checked against an eight-point standard. & That this artifact has been independently reproduced -- that item is pending, checked directly against the issue tracker, not assumed. \\
\bottomrule
\end{longtable}
\end{center}

\noindent\textbf{Proof status.} Every numbered formal claim carries a \texttt{PROOF-STATUS} tag in the companion repository's own \texttt{appendix-a-proofs.md}: \texttt{machine-checked} (an executable checker module verifies the claim over an explicitly enumerated finite domain), \texttt{checked-scope-only} (the certified scope is narrower than the general prose, stated exactly), or \texttt{pending-human-review} (asserted, not yet machine-verified) -- \texttt{proven} is never written. Every Section~\ref{sec:evidence} hypothesis (CH-B1, CH-B2, CH-C1, CH-C2) carries the same discipline inline, naming its own checker and enumerated domain rather than asserting a bare result.

\textbf{Acknowledgements.} Drafting, engineering, formal derivation, and citation verification were AI-assisted (Claude); the author is solely responsible for all claims.

\nocite{*}
\bibliographystyle{plain}
\bibliography{refs}

\end{document}
